% Theorem D''': uniform-in-time propagation of chaos in the learning regime (the averaged slow loop)
% (PARADIGM.md, section 5d).
\section{Theorem D\texorpdfstring{$'''$}{'''}: uniform-in-time propagation of chaos in the learning regime (the
averaged slow loop)}\label{sec:thmDppp}

Theorem D$''$ ends where the leak stops dominating the loop gain: in the learning regime its small-gain argument has
nothing to say, and the long-horizon run of \code{exp_chaos_uniform.py} showed the coupled distance growing over the
horizon. This section isolates the structure that the small-gain argument throws away. Under the synchronous coupling
the interaction between the particles enters through the fitted function $\fhat_\Theta$, which is linear in the
empirical measure; the loss $\bar L$ is a convex quadratic of it (flat convexity, Theorem H(i)); and the interaction
term of the coupling identity is therefore the negative of a squared $L^2(\rho)$-norm at linear order --- it
contracts, whatever the gains. What remains and can expand is the self term, whose curvature is the Hessian of the
first-variation potential, $\sigma_0\Ginv - R[\mu](\vartheta)$ with $R$ the \emph{residual-weighted} curvature of the
features --- the same object as Proposition K's $R_2$ at a collapsed configuration. Proposition K says that the
collapsed configuration is unstable when this curvature exceeds the leak \emph{with the collapsed residual}; at a
fitted configuration the residual is the approximation error, and the second-order optimality condition of the limit
measure says that the curvature is \emph{below} the leak at every atom. The theorem below turns this into a
uniform-in-time statement for the averaged slow loop at $\tau = 0$, with explicit constants, in the regime in which
the network learns.

\paragraph{Setting.} The averaged slow dynamics of Section~\ref{sec:thmJprime} at $\tau_0 = 0$ (Theorem J$'$'s
leading order $\bar L_0$ without dither; slow time $s = t/k_e$, leak $\sigma_0 = k_e\sigma$, $f = f(x)$ continuous,
$\rho$ the occupation measure of the $T_0$-periodic reference, $F_\rho := \sup_{\supp\rho}\norm{f}$, saturated
features, $K_\Phi := \sup\norm{\tilde\Phi}$), written as the $G$-preconditioned gradient flow of
$J_m(\Theta) = m\bar L(\mu_\Theta) + \frac{\sigma_0}2\sum_i\norm{\vartheta_i - \vartheta^0}^2_{\Ginv}$:
\begin{gather}
  \frac{\dd\vartheta_i}{\dd s} = -G\nabla_{\vartheta_i}J_m(\Theta) = -G\nabla_\vartheta\mathcal{V}[\mu_\Theta](\vartheta_i),\qquad
  i = 1..m,
  \tag{$\mathfrak{A}_m$}\\
  \mathcal{V}[\mu](\vartheta) := V_L[\mu](\vartheta) + \frac{\sigma_0}2\norm{\vartheta - \vartheta^0}^2_{\Ginv},\notag
\end{gather}
with $V_L[\mu](\vartheta) = -\langle r_\mu^\T\tilde\Phi(\cdot\,;\vartheta)\rangle_\rho$ and $r_\mu = f - \fhat[\mu]$,
$\fhat[\mu] := \int\tilde\Phi(\cdot\,;\vartheta')\dd\mu(\vartheta')$, as in Section~\ref{sec:thmH}. Its mean-field
equation is the Wasserstein gradient flow $\Aavg$ of Section~\ref{sec:thmH} at $\tau = 0$ in slow time,
$\partial_s\mu_s = \nabla\cdot(\mu_sG\nabla\mathcal{V}[\mu_s])$, solved by the characteristic flow
\begin{equation}
  \frac{\dd X_s(\vartheta)}{\dd s} = -G\nabla\mathcal{V}[\mu_s](X_s(\vartheta)),\qquad X_0 = \mathrm{id},\qquad
  \mu_s = (X_s)_\#\mu_0
  \tag{$\mathfrak{A}_0$}
\end{equation}
(a McKean--Vlasov equation without noise whose vector field has linear growth through the leak, is Lipschitz in
$\vartheta$ and, through $r_\mu$, Lipschitz in $\mu$ for $W_1$; existence, uniqueness and the representation by
characteristics follow from Dobrushin's argument~\cite{S35} --- the synchronous-coupling proof of Theorem D$'$(b)
with the plant removed --- and the velocity is bounded on the invariant ball of Step~0). The constants are taken in
the $\Ginv$ metric on parameters (``whitened''):
\[
  L_\vartheta^G := \sup_{x,\vartheta}\norm{\nabla_\vartheta\tilde\Phi(x;\vartheta)G^{1/2}},\qquad
  L_2^G := \sup_{x,\vartheta,\norm{u}=1}\norm{G^{1/2}\nabla^2_\vartheta\bigl(u^\T\tilde\Phi(x;\vartheta)\bigr)G^{1/2}},
\]
so that $\norm{\tilde\Phi(x;\vartheta) - \tilde\Phi(x;\vartheta')}\le L_\vartheta^G\norm{\vartheta - \vartheta'}_{\Ginv}$,
$\norm{\nabla_\vartheta\tilde\Phi(x;\vartheta)v}\le L_\vartheta^G\norm{v}_{\Ginv}$ and
$\norm{(\nabla_\vartheta\tilde\Phi(x;\vartheta) - \nabla_\vartheta\tilde\Phi(x;\vartheta'))v}
\le L_2^G\norm{\vartheta - \vartheta'}_{\Ginv}\norm{v}_{\Ginv}$ ($L_\vartheta^G\le\lambda_G^{1/2}L_\vartheta$ and
$L_2^G\le\lambda_GL_2$ in the constants of Section~\ref{sec:thmDprime};
\code{chaos_learning.whitened_lipschitz_constants} computes the suprema by search, with closed-form bounds). The
\emph{residual-weighted curvature} of $\mu$ at $\vartheta$ is the symmetric matrix
\begin{gather*}
  R[\mu](\vartheta) := \Bigl\langle\sum_kr_\mu^k(x)\nabla^2_\vartheta\tilde\Phi^k(x;\vartheta)\Bigr\rangle_\rho
  = -\nabla^2_\vartheta V_L[\mu](\vartheta),\qquad R^G := G^{1/2}RG^{1/2},\\
  \kappa_R[\mu](A) := \sup_{\vartheta\in A}\lambda_{\max}\bigl(R^G[\mu](\vartheta)\bigr),
\end{gather*}
so that $G^{1/2}\nabla^2_\vartheta\mathcal{V}[\mu](\vartheta)G^{1/2} = \sigma_0I - R^G[\mu](\vartheta)$; for
$\mu = \delta_{\bar\vartheta}$ and $\vartheta = \bar\vartheta$ it is Proposition K's $R_2$ and
$\sigma_{\mathrm{crit}}(\bar\vartheta) = \kappa_R[\delta_{\bar\vartheta}](\{\bar\vartheta\})$. It satisfies
$\norm{R^G[\mu](\vartheta)}\le L_2^G\langle\norm{r_\mu}\rangle_\rho\le L_2^G(F_\rho + K_\Phi)$ and
$\norm{R^G[\mu](\vartheta) - R^G[\nu](\vartheta)}\le L_2^G\norm{\fhat[\mu] - \fhat[\nu]}_{L^1(\rho)}$.
$N_\varrho(A) := \{\vartheta : \operatorname{dist}_{\Ginv}(\vartheta,A)\le\varrho\}$.

\smallskip
\noindent\emph{The coupling.} $\vartheta_i(0)$, $i = 1..m$, i.i.d.\ with law $\mu_0$ of compact support,
$R_0 := \sup_{\supp\mu_0}\norm{\vartheta - \vartheta^0}_{\Ginv}$; $\Theta(s)$ solves \textup{($\mathfrak{A}_m$)} and
the reference particles are $\bar\vartheta_i(s) := X_s(\vartheta_i(0))$, i.i.d.\ with law $\mu_s$ and
$\bar\vartheta_i(s)\in\supp\mu_s$. Let $\delta_i := \vartheta_i - \bar\vartheta_i$,
\[
  D(s) := \frac1m\sum_i\norm{\delta_i(s)}^2_{\Ginv},\qquad
  Y(s) := \norm{\fhat_{\bar\Theta(s)} - \fhat[\mu_s]}_{L^2(\rho)},
\]
$Y$ being the sampling fluctuation of the reference's fitted function
($\E Y(s)^2 = \frac1m\langle\operatorname{Var}_{\mu_s}\tilde\Phi(x;\cdot)\rangle_\rho\le K_\Phi^2/m$). A priori
(Step~0 below) every particle of either system stays in the ball
$\norm{\vartheta - \vartheta^0}_{\Ginv}\le R_\Theta := \max\{R_0,\ L_\vartheta^G(F_\rho + K_\Phi)/\sigma_0\}$, so
$D\le D_{\max} := 4R_\Theta^2$, and the characteristics move at speed at most
$V_{\max} := L_\vartheta^G(F_\rho + K_\Phi) + \sigma_0R_\Theta$.

\begin{thm}{Theorem}{D$'''$}[uniform-in-time propagation of chaos in the learning regime, $\tau = 0$]
Assume
\begin{itemize}[leftmargin=3em]
  \item[\textup{(H)}] \emph{(non-degenerate stability along the mean-field flow)} there are $s_1\ge0$, $\kappa > 0$
    and $\varrho' > 0$ such that
    \[
      \kappa_R[\mu_s]\bigl(N_{\varrho'}(\supp\mu_s)\bigr)\le\sigma_0 - \kappa\qquad\text{for every }s\ge s_1,
    \]
\end{itemize}
and put $a(s) := \sigma_0 - \kappa_R[\mu_s](N_{\varrho'}(\supp\mu_s))$ (so $a\ge\kappa$ on $[s_1,\infty)$ and
$a\ge\sigma_0 - L_2^G(F_\rho + K_\Phi)$ everywhere) and $\Lambda_1 := \int_0^{s_1}(\kappa - a(s))_+\dd s < \infty$.
Define
\begin{gather*}
  \varrho := \min\bigl\{(2\kappa)^{1/2}/L_2^G,\ \varrho'\kappa e^{-\Lambda_1}/(4(L_\vartheta^G)^2)\bigr\},\\
  \theta := \min\bigl\{\kappa/(8L_2^G),\ \varrho\kappa e^{-\Lambda_1}/(4L_\vartheta^G),\ \varrho'\kappa e^{-\Lambda_1}/(4L_\vartheta^G)\bigr\},\\
  \Delta := \theta/(4L_\vartheta^GV_{\max}),\qquad C_1 := 4e^{2\Lambda_1}(L_\vartheta^G)^2K_\Phi^2/\kappa^2,\qquad
  m_0 := 16K_\Phi^2/\theta^2.
\end{gather*}
Then for every $m\ge m_0$:
\begin{enumerate}[label=\textup{(\alph*)}]
  \item \emph{(exponentially long horizons)} for every $S > 0$ and every $s\in[0,S]$,
    \[
      \E D(s)\le C_1/m + 4R_\Theta^2(S/\Delta + 2)\exp\bigl(-m\theta^2/(32K_\Phi^2)\bigr);
    \]
    in particular $\sup_{s\le S}\E D(s)\le(C_1 + 12R_\Theta^2)/m$ for all $S\le\Delta e^{m\theta^2/(64K_\Phi^2)}$ and
    $m\ge m_0$ with $e^{m\theta^2/(64K_\Phi^2)}\ge m$.
  \item \emph{(uniform in time)} if moreover \textup{(H$'$)} the characteristic flow converges exponentially on the
    support of $\mu_0$, $\sup_{\vartheta\in\supp\mu_0}\norm{X_s(\vartheta) - X_\infty(\vartheta)}_{\Ginv}\le C_2e^{-\kappa_2s}$
    for all $s\ge0$, then for every $s\ge0$
    \[
      \E D(s)\le C_1/m + 4R_\Theta^2(R_{\mathrm{eff}}/\Delta + 3)\exp\bigl(-m\theta^2/(32K_\Phi^2)\bigr),\qquad
      R_{\mathrm{eff}} := \kappa_2^{-1}\bigl(\log(8L_\vartheta^GC_2/\theta)\bigr)_+.
    \]
  \item \emph{(the learned function and the particles)} in either case, with the same right-hand side $B(s)$ for
    $\E D(s)$,
    \begin{gather*}
      \E\norm{\fhat_{\Theta(s)} - \fhat[\mu_s]}^2_{L^2(\rho)}\le2(L_\vartheta^G)^2B(s) + 2K_\Phi^2/m,\\
      \E W_2(\mu_{\Theta(s)},\mu_s)\le B(s)^{1/2} + \E W_2(\mu_{\bar\Theta(s)},\mu_s),
    \end{gather*}
    and on the event $\mathcal{G}_S := \{\sup_{s\le S}Y(s)\le\theta\}$ (case \textup{(a)}; $\mathcal{G}_\infty$ in
    case \textup{(b)}), whose complement is the tail above, every particle stays within $3\varrho'/4$ of its reference
    up to time $S$: $\max_i\sup_{s\le S}\norm{\delta_i(s)}_{\Ginv}\le3\varrho'/4$ --- no particle changes basin.
\end{enumerate}
The constants depend on $(\kappa,\varrho',\Lambda_1,L_\vartheta^G,L_2^G,K_\Phi,F_\rho,\sigma_0,R_0)$ and not on $m$,
$s$ or $S$.
\end{thm}
\status{Proved (synchronous coupling with the three-term identity of Step~1, a bootstrap on the mean square and the
per-particle distance, and McDiarmid's inequality~\cite{S36} for the reference's sampling fluctuation; the
finite-horizon part of the transient is Theorem D$'$'s mechanism and appears in $\Lambda_1$). Hypothesis (H) is a
property of the deterministic mean-field flow and can be computed; it is the strict second-order optimality condition
of the limit configuration with a margin and a neighbourhood (Remark D$'''$.1), and it is compatible with the learning
regime $\sigma_0 < \sigma_{\mathrm{crit}}$ of Proposition K. It requires the support of $\mu_s$ to stay away from the
ridges of $\mathcal{V}[\mu_s]$ for $s\ge s_1$: for an initial cloud whose support is a union of pieces each flowing
into one basin (a warm start, or a cluster per bump) this holds; for a connected cold start whose mass splits between
several basins it fails for every $s$ (Remark D$'''$.2), and the numerics below show what the coupling does there.
Numerics: \code{exp_chaos_learning.py} $\to$ \code{results/chaos/learning_table.md} (below). The theorem is about the
averaged slow dynamics; its transfer to the real loop at finite $k_e$ is open (Remark D$'''$.4). Its constants are
worst-case and its regime is asymptotic: on the instance below $m_0$ is of order $10^{23}$ (Remark D$'''$.6), and the
finite-$m$ behaviour is described sharply by Proposition D$'''$.5. Referee round 8
(\code{review/REVIEW_round8_Dppp.md}) re-derived every constant of Steps 0--9, found no counterexample to (2.1),
(3.1), (4.1) on an adversarial search, and corrected the wording of Remarks D$'''$.1--D$'''$.2, the good event per
horizon and the well-posedness sentence (all applied). Referee round 9 (\code{review/REVIEW_round9_D3b.md}) reviewed
Proposition D$'''$.5, Corollary D$'''$.7 and the numerics: all hold with fixes, the main one being that the finite-$m$
replicated configuration is a minimiser only for admissible weights (Proposition D$'''$.5(iv)) --- one of the four
$m = 16$ warm seeds sits on a saddle, which the numerics paragraph now says. Novelty (literature audit): the technique (synchronous
coupling + bootstrap + bounded differences) is standard; the sign structure of the interaction and the local,
residual-weighted hypothesis --- and with them a uniform-in-time statement in a regime where the network learns ---
are new in propagation of chaos for adaptive control; the theorem's regime is asymptotic in $m$ and the finite-$m$ law
is a delta-method statement, new in its object.}

\paragraph{Literature audit} (\code{LITERATURE_TRANSPORT.md}, D$'''$).
\begin{itemize}
  \item \emph{Prior.} Uniform-in-time propagation of chaos by synchronous coupling under \emph{global} convexity of
    the confinement and convexity of the interaction is Malrieu's~\cite{S37}; non-convex confinements with
    \emph{small} Lipschitz interactions are handled by reflection coupling~\cite{S38}; flat-convex energies with dither
    by the entropy method~\cite{S5}. The Hessian decomposition ``Gram + residual-weighted feature curvature'' is the
    standard tool of the interpolation\,/\,PL-condition literature~\cite{S39} and of lazy-training analyses, where the
    second term vanishes at interpolation.
  \item \emph{Distinction.} (i) The interaction of the fitted network is neither small nor convex in the particle
    coordinates; its sign comes from the flat convexity of $\bar L$ ($\fhat$ linear in the measure), which makes
    $-\norm{\fhat_\Theta - \fhat_{\bar\Theta}}^2$ the leading interaction term under the coupling --- the finite-$m$,
    two-measure version of \cite{S39}'s decomposition. (ii) The convexity needed is local:
    $\sigma_0I - G^{1/2}R[\mu_s]G^{1/2}\succeq\kappa I$ on a neighbourhood of the support of the mean-field flow,
    i.e.\ the strict second-order optimality condition of the fitted configuration, evaluated at the actual residual
    --- not a global convexity (which would force the collapsed configuration to be stable and the network not to
    learn, Theorem D$''$). (iii) The window $\kappa_R[\mu_\infty] < \sigma_0 < \sigma_{\mathrm{crit}}$ between the
    collapsed instability (Proposition K) and the fitted stability is a design statement about the leak that has no
    counterpart in the cited work. (iv) The limits are stated with the result: a connected initial support that splits
    between basins violates the hypothesis for all times (the characteristic flow is a homeomorphism), and the
    coupling then has a floor of pairs in different basins; the real loop at finite $k_e$ is not covered.
  \item \emph{Proposition D$'''$.5 (the weight-sensitivity law).} At finite $m$ the stationary coupled distance is the
    multinomial fluctuation of the cluster weights propagated through the implicit function $w\mapsto\Theta^*(w)$: a
    delta-method statement (the Gaussian limit of a smooth function of a multinomial), standard in statistics; what is
    specific is its object (the shared-particle coupling of a particle approximation to a stationary $K$-atom measure,
    with the covariance $(1/m - 1/M)(\operatorname{diag}\bar w - \bar w\bar w^\T)$) and its role as the sharp form of
    the theorem's constant (stiffness$^{-1}$ $\times$ feature correlations against $\kappa^{-1}L_\vartheta^GK_\Phi$).
\end{itemize}

\begin{proof}
\emph{Step 0 (a priori bounds).} For either system,
$\frac{\dd}{\dd s}\norm{\vartheta - \vartheta^0}^2_{\Ginv}
= 2\ip{\vartheta - \vartheta^0}{\langle\nabla_\vartheta(r^\T\tilde\Phi(\cdot\,;\vartheta))\rangle_\rho}
- 2\sigma_0\norm{\vartheta - \vartheta^0}^2_{\Ginv}$ with
$\abs{\ip{\vartheta - \vartheta^0}{\langle\nabla_\vartheta(r^\T\tilde\Phi)\rangle_\rho}}
= \abs{\langle r^\T\nabla_\vartheta\tilde\Phi(\vartheta - \vartheta^0)\rangle_\rho}
\le\langle\norm{r}\rangle_\rho L_\vartheta^G\norm{\vartheta - \vartheta^0}_{\Ginv}
\le(F_\rho + K_\Phi)L_\vartheta^G\norm{\vartheta - \vartheta^0}_{\Ginv}$, so
$D^+\norm{\vartheta - \vartheta^0}_{\Ginv}\le L_\vartheta^G(F_\rho + K_\Phi) - \sigma_0\norm{\vartheta - \vartheta^0}_{\Ginv}$
and $\norm{\vartheta(s) - \vartheta^0}_{\Ginv}\le R_\Theta$; hence $\norm{\delta_i}_{\Ginv}\le2R_\Theta$,
$D\le D_{\max}$, and
$\norm{\dd\vartheta/\dd s}_{\Ginv} = \norm{G^{1/2}\langle\nabla_\vartheta(r^\T\tilde\Phi)\rangle_\rho - \sigma_0G^{-1/2}(\vartheta - \vartheta^0)}\le V_{\max}$.
Two Lipschitz facts: for $\Theta$, $\bar\Theta$ with rows $\vartheta_j$, $\bar\vartheta_j$ and
$A_j(x) := \int_0^1\nabla_\vartheta\tilde\Phi(x;\bar\vartheta_j + t\delta_j)\dd t$,
\begin{gather}
  \fhat_\Theta - \fhat_{\bar\Theta} = \frac1m\sum_jA_j\delta_j =: h,\qquad
  \norm{h}_{L^2(\rho)}\le L_\vartheta^G\frac1m\sum_j\norm{\delta_j}_{\Ginv}\le L_\vartheta^GD^{1/2},
  \tag{0.1}\\
  \begin{split}
  &\norm{R^G[\mu_{\bar\Theta}](\vartheta) - R^G[\mu_s](\vartheta)}\le L_2^G\norm{r_{\bar\Theta} - r_{\mu_s}}_{L^1(\rho)}\le L_2^GY,\\
  &\norm{(\nabla_\vartheta\tilde\Phi(x;\vartheta_i) - A_i(x))v}\le\frac{L_2^G}2\norm{\delta_i}_{\Ginv}\norm{v}_{\Ginv}.
  \end{split}
  \tag{0.2}
\end{gather}

\smallskip
\emph{Step 1 (the identity).} Since $\dd\delta_i/\dd s = -G[\nabla\mathcal{V}[\mu_\Theta](\vartheta_i)
- \nabla\mathcal{V}[\mu_s](\bar\vartheta_i)]$ and $\ip{\delta}{Gv}_{\Ginv} = \ip{\delta}{v}$,
\begin{gather*}
  \tfrac12\frac{\dd D}{\dd s} = -\frac1m\sum_i\ip{\delta_i}{\nabla\mathcal{V}[\mu_\Theta](\vartheta_i) - \nabla\mathcal{V}[\mu_s](\bar\vartheta_i)}
  = T_A + T_B + T_C,\\
  T_A := -\frac1m\sum_i\ip{\delta_i}{\nabla\mathcal{V}[\mu_{\bar\Theta}](\vartheta_i) - \nabla\mathcal{V}[\mu_{\bar\Theta}](\bar\vartheta_i)},\\
  T_B := -\frac1m\sum_i\ip{\delta_i}{\nabla V_L[\mu_\Theta](\vartheta_i) - \nabla V_L[\mu_{\bar\Theta}](\vartheta_i)},\\
  T_C := -\frac1m\sum_i\ip{\delta_i}{\nabla V_L[\mu_{\bar\Theta}](\bar\vartheta_i) - \nabla V_L[\mu_s](\bar\vartheta_i)}
\end{gather*}
(the leak is linear and sits in $T_A$; the three brackets telescope to
$\nabla\mathcal{V}[\mu_\Theta](\vartheta_i) - \nabla\mathcal{V}[\mu_s](\bar\vartheta_i)$). The self term $T_A$
compares the \emph{reference's} potential at the two points; the interaction term $T_B$ compares the two empirical
measures at the $m$-system's point; $T_C$ is the reference's sampling error.
\code{chaos_learning.AveragedSlowLoop.decomposition} evaluates the three terms and checks the identity to rounding.

\smallskip
\emph{Step 2 (the self term is local convexity).}
$\nabla\mathcal{V}[\mu_{\bar\Theta}](\vartheta_i) - \nabla\mathcal{V}[\mu_{\bar\Theta}](\bar\vartheta_i)
= \int_0^1\nabla^2\mathcal{V}[\mu_{\bar\Theta}](\bar\vartheta_i + t\delta_i)\dd t\cdot\delta_i
= \sigma_0\Ginv\delta_i - \bar R_i\delta_i$ with $\bar R_i := \int_0^1R[\mu_{\bar\Theta}](\bar\vartheta_i + t\delta_i)\dd t$,
so $T_A = -\sigma_0D + \frac1m\sum_i\ip{\delta_i}{\bar R_i\delta_i}$ and
$\ip{\delta}{R\delta}\le\lambda_{\max}(R^G)\norm{\delta}^2_{\Ginv}$. If $\norm{\delta_i}_{\Ginv}\le\varrho'$ the
segment lies in $N_{\varrho'}(\supp\mu_s)$ ($\bar\vartheta_i\in\supp\mu_s$), and by (0.2)
$\lambda_{\max}(R^G[\mu_{\bar\Theta}])\le\kappa_R[\mu_s](N_{\varrho'}(\supp\mu_s)) + L_2^GY$ there. Hence, as long
as $\max_i\norm{\delta_i}_{\Ginv}\le\varrho'$,
\begin{equation}
  T_A\le-\bigl(a(s) - L_2^GY(s)\bigr)D.
  \tag{2.1}
\end{equation}

\smallskip
\emph{Step 3 (the interaction term is a negative square at linear order).}
\[
  \nabla V_L[\mu_\Theta](\vartheta_i) - \nabla V_L[\mu_{\bar\Theta}](\vartheta_i)
  = \langle\nabla_\vartheta\tilde\Phi(\cdot\,;\vartheta_i)^\T(\fhat_\Theta - \fhat_{\bar\Theta})\rangle_\rho
  = \langle\nabla_\vartheta\tilde\Phi(\cdot\,;\vartheta_i)^\T h\rangle_\rho,
\]
so, with the $L^2(\rho;\R^n)$ inner product $\langle\langle\cdot,\cdot\rangle\rangle$,
\begin{equation}
  \begin{split}
  T_B &= -\Bigl\langle\Bigl\langle\frac1m\sum_i\nabla_\vartheta\tilde\Phi(\cdot\,;\vartheta_i)\delta_i,\,h\Bigr\rangle\Bigr\rangle
  = -\norm{h}^2 - \Bigl\langle\Bigl\langle\frac1m\sum_i(\nabla_\vartheta\tilde\Phi(\cdot\,;\vartheta_i) - A_i)\delta_i,\,h\Bigr\rangle\Bigr\rangle\\
  &\le-\norm{h}^2 + \norm{h}\frac{L_2^G}2D\le\frac{(L_2^G)^2D^2}{16}
  \end{split}
  \tag{3.1}
\end{equation}
by (0.2) and $\max_x(-x^2 + cx) = c^2/4$. The leading part
$-\norm{h}^2 = -\norm{\fhat_\Theta - \fhat_{\bar\Theta}}^2_{L^2(\rho)}$ is the flat convexity of $\bar L$: the
interaction contracts, and only the Taylor remainder of the features, quadratic in $D$, is charged. This is the step
that Theorem D$''$'s small-gain bound replaces by $+L_\vartheta^2D$-type terms.

\smallskip
\emph{Step 4 (the sampling term).}
$\nabla V_L[\mu_{\bar\Theta}](\bar\vartheta_i) - \nabla V_L[\mu_s](\bar\vartheta_i)
= \langle\nabla_\vartheta\tilde\Phi(\cdot\,;\bar\vartheta_i)^\T(\fhat_{\bar\Theta} - \fhat[\mu_s])\rangle_\rho$, so
$\abs{\ip{\delta_i}{\cdot}}\le L_\vartheta^G\norm{\delta_i}_{\Ginv}Y$ and
\begin{equation}
  T_C\le L_\vartheta^GD^{1/2}Y.
  \tag{4.1}
\end{equation}

\smallskip
\emph{Step 5 (the differential inequalities).} While $\max_i\norm{\delta_i}_{\Ginv}\le\varrho'$,
\begin{equation}
  \tfrac12\frac{\dd D}{\dd s}\le-(a - L_2^GY)D + \frac{(L_2^G)^2D^2}{16} + L_\vartheta^GD^{1/2}Y,
  \tag{5.1}
\end{equation}
and per particle, from the same three brackets with
$\abs{T_{B,i}}\le\norm{\delta_i}L_\vartheta^G\norm{h}\le\norm{\delta_i}(L_\vartheta^G)^2D^{1/2}$ by (0.1) and
$\abs{T_{C,i}}\le\norm{\delta_i}L_\vartheta^GY$,
\begin{equation}
  D^+\norm{\delta_i}_{\Ginv}\le-(a - L_2^GY)\norm{\delta_i}_{\Ginv} + (L_\vartheta^G)^2D^{1/2} + L_\vartheta^GY.
  \tag{5.2}
\end{equation}

\smallskip
\emph{Step 6 (the bootstrap on the good event).} Fix a horizon $S\le\infty$, let
$\mathcal{G}_S := \{\sup_{s\le S}Y(s)\le\theta\}$ and
$T^* := \inf\{s\le S : D(s)\ge\varrho^2\text{ or }\max_i\norm{\delta_i(s)}_{\Ginv}\ge\varrho'\}$ ($T^* > 0$ since
$\delta(0) = 0$; everything below is on $[0,S]$). On $\mathcal{G}_S$ and $[0,T^*)$: $L_2^GY\le\kappa/8$,
$(L_2^G)^2D/16\le(L_2^G)^2\varrho^2/16\le\kappa/8$, and
$L_\vartheta^GD^{1/2}Y\le(\kappa/8)D + (2(L_\vartheta^G)^2/\kappa)Y^2$ (Young), so (5.1) gives
$\dd D/\dd s\le-2(a - 3\kappa/8)D + (4(L_\vartheta^G)^2/\kappa)Y^2$. Since
$a - 3\kappa/8\ge\kappa/2 - (\kappa - a)_+$, $\int_r^s(a - 3\kappa/8)\ge(\kappa/2)(s - r) - \Lambda_1$ for $r\le s$,
and Gr\"onwall gives
\begin{equation}
  \begin{split}
  D(s)\le Z(s) &:= 4e^{2\Lambda_1}(L_\vartheta^G)^2\kappa^{-1}\int_0^se^{-\kappa(s-r)}Y(r)^2\dd r\\
  &\le4e^{2\Lambda_1}(L_\vartheta^G)^2\kappa^{-2}\theta^2\le\varrho^2/4\qquad\text{on }\mathcal{G}_S\cap[0,T^*)
  \end{split}
  \tag{6.1}
\end{equation}
by the second constraint on $\theta$. Per particle, (5.2) with $D^{1/2}\le\varrho/2$ and $Y\le\theta$, and
$\int_r^s(a - \kappa/8)\ge(\kappa/2)(s - r) - \Lambda_1$, gives
$\norm{\delta_i(s)}_{\Ginv}\le2e^{\Lambda_1}\kappa^{-1}[(L_\vartheta^G)^2\varrho/2 + L_\vartheta^G\theta]
\le\varrho'/4 + \varrho'/2 = 3\varrho'/4$ by the second constraint on $\varrho$ and the third on $\theta$. Both
thresholds are therefore strict on $[0,T^*)$, so $T^* = S$ on $\mathcal{G}_S$ by continuity: $D\le Z$ and
$\max_i\norm{\delta_i}\le3\varrho'/4$ on $[0,S]$ on $\mathcal{G}_S$. $Z$ does not involve the $m$-system.

\smallskip
\emph{Step 7 (expectation).} For $s\le S$,
\begin{multline*}
  \E D(s)\le\E[Z(s)] + D_{\max}\mathbb{P}(\mathcal{G}_S^c)
  \le4e^{2\Lambda_1}(L_\vartheta^G)^2\kappa^{-1}\int_0^se^{-\kappa(s-r)}\E Y(r)^2\dd r + D_{\max}\mathbb{P}(\mathcal{G}_S^c)\\
  \le C_1/m + 4R_\Theta^2\mathbb{P}(\mathcal{G}_S^c),
\end{multline*}
using $\E Y(r)^2\le K_\Phi^2/m$.

\smallskip
\emph{Step 8 (the good event).} $Y(s)$ is a function of the i.i.d.\ sample $(\vartheta_j(0))_{j\le m}$:
$Y(s) = \norm{\frac1m\sum_j[\tilde\Phi(\cdot\,;X_s(\vartheta_j(0))) - \int\tilde\Phi(\cdot\,;X_s(\vartheta))\dd\mu_0]}_{L^2(\rho)}$.
Changing one $\vartheta_j(0)$ changes it by at most $2K_\Phi/m$ (bounded differences), and
$\E Y(s)\le(\E Y^2)^{1/2}\le K_\Phi/m^{1/2}$; McDiarmid's inequality~\cite{S36} gives
$\mathbb{P}(Y(s)\ge K_\Phi/m^{1/2} + t)\le\exp(-mt^2/(2K_\Phi^2))$, so for $m\ge m_0 = 16K_\Phi^2/\theta^2$
($K_\Phi/m^{1/2}\le\theta/4$, $t = \theta/4$)
\begin{equation}
  \mathbb{P}(Y(s)\ge\theta/2)\le\exp\bigl(-m\theta^2/(32K_\Phi^2)\bigr)\qquad\text{for every fixed }s.
  \tag{8.1}
\end{equation}
In time, $\abs{Y(s) - Y(s')}\le2L_\vartheta^GV_{\max}\abs{s - s'}$ (each $X_s(\vartheta)$ moves at speed
$\le V_{\max}$ and $\tilde\Phi$ is $L_\vartheta^G$-Lipschitz, for the sample and for the mean), so on a grid of step
$\Delta = \theta/(4L_\vartheta^GV_{\max})$, $Y\le\theta$ on $[s_k,s_k + \Delta]$ whenever $Y(s_k)\le\theta/2$. (a) The
grid $\{k\Delta : k\Delta\le S\}\cup\{S\}$ has at most $S/\Delta + 2$ points, and a union bound with (8.1) gives
$\mathbb{P}(\mathcal{G}_S^c) = \mathbb{P}(\sup_{s\le S}Y > \theta)\le(S/\Delta + 2)\exp(-m\theta^2/(32K_\Phi^2))$;
with Step~7 this is (a). For $S\le\Delta e^{m\theta^2/(64K_\Phi^2)}$,
$(S/\Delta + 2)e^{-m\theta^2/(32K_\Phi^2)}\le3e^{-m\theta^2/(64K_\Phi^2)}\le3/m$ when
$e^{m\theta^2/(64K_\Phi^2)}\ge m$. (b) Under (H$'$), for $s\ge R_{\mathrm{eff}}$ every $X_s(\vartheta)$ is within
$\theta/(8L_\vartheta^G)$ of $X_\infty(\vartheta)$ on $\supp\mu_0$, so $Y(s)\le Y(\infty) + \theta/4$ with
$Y(\infty) := \norm{\frac1m\sum_j[\tilde\Phi(\cdot\,;X_\infty(\vartheta_j(0))) - \int\tilde\Phi(\cdot\,;X_\infty)\dd\mu_0]}$,
another bounded-differences statistic with $\E Y(\infty)\le K_\Phi/m^{1/2}$; the grid on $[0,R_{\mathrm{eff}}]$ plus
the point at infinity has at most $R_{\mathrm{eff}}/\Delta + 3$ points, and
$\mathbb{P}(\mathcal{G}_\infty^c) = \mathbb{P}(\sup_{s\ge0}Y > \theta)\le(R_{\mathrm{eff}}/\Delta + 3)\exp(-m\theta^2/(32K_\Phi^2))$.

\smallskip
\emph{Step 9 (the learned function and the measures).}
$\norm{\fhat_\Theta - \fhat[\mu_s]}\le\norm{h} + Y\le L_\vartheta^GD^{1/2} + Y$ by (0.1), so
$\E\norm{\cdot}^2\le2(L_\vartheta^G)^2\E D + 2K_\Phi^2/m$. The coupling $(\vartheta_i,\bar\vartheta_i)$ is a coupling
of $\mu_\Theta$ and $\mu_{\bar\Theta}$, so $W_2^2(\mu_\Theta,\mu_{\bar\Theta})\le D$ ($\Ginv$ metric), and the
triangle inequality gives (c). The per-particle statement is Step~6.
\end{proof}

\begin{rmk}{Remark}{D$'''$.1}[what the hypothesis says, and the window]
(H) at $s = \infty$ reads: on a neighbourhood of every atom of the limit measure,
$G^{1/2}\nabla^2\mathcal{V}[\mu_\infty]G^{1/2} = \sigma_0I - R^G[\mu_\infty]\succeq\kappa I$ --- each atom is a
non-degenerate local minimiser of the first-variation potential $\mathcal{V}[\mu_\infty]$, with margin $\kappa$: the
strict form of the second-order optimality condition for measures (the local Chizat--Bach condition, as in
Proposition K(ii)). Its non-strict form, $\sigma_0\Ginv - R[\mu_\infty](\vartheta^*)\succeq0$ at every atom
$\vartheta^*$ of positive mass, is necessary for an atomic stationary measure to be a local minimiser of $\mathcal{G}$
in the Wasserstein sense: splitting an atom of mass $w$ along a direction of negative curvature of
$\mathcal{V}[\mu_\infty]$ lowers $\mathcal{G}$ at second order (Proposition K(i)--(ii) with the atom's own residual).
Along the flow, (H) asks for this with the \emph{current} residual $r_{\mu_s}$ from the time $s_1$ on. It does not
mention the collapsed configuration: $\sigma_{\mathrm{crit}} = \kappa_R[\delta_{\bar\vartheta}](\{\bar\vartheta\})$
is the same curvature \emph{with the collapsed residual}, which is the residual of the best single kernel and is
large when the plant needs several. The learning regime $\sigma_0 < \sigma_{\mathrm{crit}}$ (the kernels separate)
and the uniform regime (H) ($\kappa_R[\mu_\infty]\le\sigma_0 - \kappa$ at the fitted configuration) therefore coexist
exactly when the fitted configuration is non-degenerately stable while the collapsed one is not --- the window
\[
  \kappa_R[\mu_\infty]\bigl(N_{\varrho'}(\supp\mu_\infty)\bigr) < \sigma_0 < \sigma_{\mathrm{crit}}(\bar\vartheta)
\]
of leaks in which the network learns \emph{and} propagation of chaos is uniform in time. Theorem D$''$'s condition,
by contrast, asks the leak to dominate the worst-case residual curvature $\lambda_GL_2(F_\rho + K_\Phi)$ \emph{and}
the interaction's Lipschitz constant $L_\vartheta^2$-type terms, which forces $\sigma_0\ge\sigma_{\mathrm{crit}}$.
The two ingredients that remove the horizon here are (i) the sign of the interaction, which is the flat convexity of
$\bar L$ (Step~3), and (ii) the residual-weighted curvature evaluated at the actual residual rather than at the worst
case (Step~2). Both are properties of the lifted linear parametrisation $\fhat = \int\tilde\Phi\dd\mu$, not of the
particular features.
\end{rmk}

\begin{rmk}{Remark}{D$'''$.2}[the transient, the support, and the floor of the coupling]
Two things limit the theorem, and both are visible in the numerics. (i) The constant carries $e^{2\Lambda_1}$, the
finite-horizon amplification of Theorem D$'$ over the time the flow needs to enter the non-degenerate regime: for a
warm start $\Lambda_1 = 0$ and $C_1 = 4(L_\vartheta^G)^2K_\Phi^2/\kappa^2$; for a cold start the expanding phase
($a(s) < \kappa$ while the residual is large) lasts a time of order $1/\sigma_0$ and $\Lambda_1$ is of order
$\sigma_0^{-1}\cdot\sup(\kappa_R - \sigma_0)$ --- the bound is uniform but quantitatively void, as D$'$'s $e^{\Lambda T}$
is. (ii) (H) is a statement about the $\varrho'$-neighbourhood of the \emph{support} of $\mu_s$. The characteristic
flow is a homeomorphism, so a connected support stays connected: if the mass of a connected cold start splits between
several basins of $\mathcal{V}[\mu_\infty]$, the support of $\mu_s$ threads through the ridges between them for every
$s$, and the filament's points converge to the saddles of $\mathcal{V}[\mu_\infty]$ on those ridges, where
$\mathcal{V}$ has a direction of negative curvature and $\kappa_R > \sigma_0$ --- (H) fails for every $s_1$
(asymptotically in $s$; pointwise the filament may well lie in a locally convex region for a while, referee round 8,
D3-11), however little mass the filament carries. The theorem then says nothing, and the coupling can indeed carry a
floor: a particle whose reference lies near a separatrix can be pushed to the other side by the $O(m^{-1/2})$
difference between the two drifts, after which the pair sits in different basins for ever, at a distance of the order
of the basin separation. A fraction $p_m$ of such pairs contributes $p_m\cdot O(R_\Theta^2)$ to $D$, so $D$ follows
the law of $p_m$ (the $\mu_0$-mass of a boundary layer of width $O(m^{-1/2})$ around the separatrices' preimage, hence
$p_m = O(m^{-1/2})$ when $\mu_0$ has a bounded density there), not $1/m$; the learned function is less sensitive,
$\norm{\fhat_\Theta - \fhat_{\bar\Theta}}\le L_\vartheta^G\frac1m\sum\norm{\delta_i}_{\Ginv}
\le L_\vartheta^G(p_m\cdot2R_\Theta + D_{\mathrm{good}}^{1/2})$, which is still $O(m^{-1/2})$. The measure-weighted
form of (H) --- a bad set of small $\mu_s$-mass, with the per-particle bootstrap run on the good particles only and
the interaction split as in Step~3 --- is the natural extension and is listed in Section~\ref{sec:program}; what the
present theorem shows is where the clean statement holds (supports inside basins) and why the cold start of Theorem
D$''$'s experiment was outside it on two counts (the transient and the connected support). (iii) The floor has a
second source at small $m$, on a warm start: when the realised cluster weights fall outside the admissible set of
Proposition D$'''$.5(iv) the replicated configuration is a saddle of $J_m$, the $m$-system lingers on it and then
splits a cluster, so that pairs end in different basins although the mean-field flow satisfies (H); its probability
is exponentially small in $m$ but appreciable at $m = 16$ on the instance below.
\end{rmk}

\begin{rmk}{Remark}{D$'''$.3}[dither]
At $\tau_0 > 0$ the synchronous coupling cancels the noise and the identity of Step~1 holds pathwise with the same
three terms, but the reference law has full support and the particles visit the ridges of $\mathcal{V}$ with the
Gibbs weight of the ridge, so (H) fails in the sense of (ii) above and the coupling carries a floor of that order.
Uniform-in-time propagation of chaos at $\tau_0 > 0$ is the province of the entropy method: with dither,
\textup{($\mathfrak{A}_m$)} is the $m$-particle mean-field Langevin dynamics of the free energy $F_\tau$ of Theorem H
in the $\Ginv$ metric (Corollary J$'$.3 is its stationary statement), and \cite{S5} proves uniform-in-time
propagation of chaos for mean-field Langevin dynamics under flat convexity, regularity of the first variation and a
uniform log-Sobolev inequality for the proximal Gibbs measures --- the three hypotheses that Theorem H(i)--(ii)
verifies for $F_\tau$, with the constant $\rho_\tau$ of H(ii). That result is imported here, not re-derived; its
constants are those of Theorem H and are exponentially small in $1/\tau_0$, where the constants of Theorem D$'''$
are those of the leak and the residual curvature. The leak-dominated regime is covered for the real loop and every
$\tau\ge0$ by Theorem D$''$.
\end{rmk}

\begin{rmk}{Remark}{D$'''$.4}[the real loop]
Theorem D$'''$ is a statement about the averaged slow dynamics \textup{($\mathfrak{A}_m$)}. Theorems J and J$'$
identify \textup{($\mathfrak{A}_m$)} as the leading order of the real loop's slow dynamics, with the corrections
$\varepsilon\cdot b$ (Theorem J$'$'s $b$, not a gradient: Proposition J$'$.2) and the fast error eliminated at
$O(\varepsilon^2)$; the transfer of a uniform-in-time statement through these corrections --- a near-identity change
of variables that is uniform in time --- is the open item of Section~\ref{sec:program}, branch 1; on a warm start it
is settled by Corollary D$'''$.7 below through Theorem J. What the theorem settles is the mechanism: in the learning
regime the uniformity rests on the fitted configuration's non-degenerate stability and on the flat convexity of the
loss, not on a dominant leak; the real-loop run on a warm start (numerics, table 5) tracks the averaged slow
dynamics over the horizon run, at the same $k_e$ at which Theorem D$''$'s cold start grew.
\end{rmk}

\begin{thm}{Proposition}{D$'''$.5}[the stationary coupled distance at finite $m$: the weight-sensitivity law]
Let $\mu^* = \sum_{k=1}^K\bar w_k\delta_{\vartheta^*_k}$ be a $K$-atom stationary measure of
\textup{($\mathfrak{A}_0$)} (each $\vartheta^*_k$ a critical point of $\mathcal{V}[\mu^*]$), and for weights
$w\in(0,\infty)^K$ let $\mathcal{G}_w(\Theta) := \frac12\langle\norm{f - \sum_kw_k\tilde\Phi(\cdot\,;\vartheta_k)}^2\rangle_\rho
+ \frac{\sigma_0}2\sum_kw_k\norm{\vartheta_k - \vartheta^0}^2_{\Ginv}$ be the mean-field energy of
$\mu_w := \sum_kw_k\delta_{\vartheta_k}$, so that $\nabla_{\vartheta_k}\mathcal{G}_w(\Theta) = w_k\nabla_\vartheta\mathcal{V}[\mu_w](\vartheta_k)$.
Assume the $K$-atom Hessian $\nabla^2_\Theta\mathcal{G}_{\bar w}(\Theta^*)$ is positive definite. Then:
\begin{enumerate}[label=\textup{(\roman*)}]
  \item \emph{(the weighted stationary point and its replication)} there is a smooth map $w\mapsto\Theta^*(w)$ on a
    neighbourhood of $\bar w$ with $\Theta^*(\bar w) = \Theta^*$ and
    $\nabla_\vartheta\mathcal{V}[\mu_{\Theta^*(w)}](\vartheta^*_k(w)) = 0$ for every $k$ (implicit function theorem);
    $S(w) := \partial\Theta^*/\partial w$ ($S_k := \partial\vartheta^*_k/\partial w\in\R^{p\times K}$) is given by
    $S = -[\nabla^2_\Theta\mathcal{G}_w]^{-1}\partial_w\nabla_\Theta\mathcal{G}_w$ at $(\Theta^*(w),w)$. For $m$
    particles with $n_k$ of them in cluster $k$ ($w^m_k := n_k/m$), the configuration $\Theta^*_m(w^m)$ with $n_k$
    copies of $\vartheta^*_k(w^m)$ is a critical point of $J_m$ ($\mu_\Theta = \mu_{w^m}$ there). Its whitened
    Hessian splits (Proposition K(i), applied cluster by cluster) into the separation blocks
    $\sigma_0I - R^G[\mu_{w^m}](\vartheta^*_k(w^m))$, each with multiplicity $n_k - 1$, and the collective block
    $\operatorname{diag}(w^m)^{-1/2}\cdot G_K^{1/2}\nabla^2_\Theta\mathcal{G}_{w^m}G_K^{1/2}\cdot\operatorname{diag}(w^m)^{-1/2}$
    (the $K$-atom Hessian, whitened and weighted). Hence $\Theta^*_m(w^m)$ is a strict local minimiser of $J_m$
    exactly when $w^m$ lies in the \emph{admissible set}
    \[
      W_{\mathrm{adm}} := \bigl\{w : \sigma_0\Ginv - R[\mu_{\Theta^*(w)}](\vartheta^*_k(w))\succ0\text{ for every }k
      \text{ and }\nabla^2\mathcal{G}_w(\Theta^*(w))\succ0\bigr\},
    \]
    a neighbourhood of $\bar w$ (both conditions are open and hold at $\bar w$ by (H) and the assumption), and a
    saddle when the splitting condition fails at some cluster: a cluster can then split, and the $m$-system leaves
    the configuration by migrating particles between basins.
  \item \emph{(the stationary distance)} Under the shared-particle coupling of Theorem D$'''$ with i.i.d.\ cluster
    labels of law $\bar w$ (the $m$ particles are the first $m$ of the $M$ reference particles), suppose
    $w^m, w^M\in W_{\mathrm{adm}}$ and that the $m$-system converges to $\Theta^*_m(w^m)$ and the reference to
    $\Theta^*_M(w^M)$. Then the coupling distance converges to
    \[
      D_\infty = \sum_kw^m_k\norm{\vartheta^*_k(w^m) - \vartheta^*_k(w^M)}^2_{\Ginv}
      = \sum_kw^m_k\norm{S_k(w^M)(w^m - w^M)}^2_{\Ginv} + O(\norm{w^m - w^M}^3),
    \]
    and, with $\tilde S_k := G^{-1/2}S_k(\bar w)$ (replacing $S(w^M)$ by $S(\bar w)$ costs a relative
    $O(\norm{w^M - \bar w})$) and $\E[(w^m - w^M)(w^m - w^M)^\T] = (1/m - 1/M)\tilde C$,
    $\tilde C := \operatorname{diag}(\bar w) - \bar w\bar w^\T$ (multinomial counts of the shared labels),
    conditionally on the admissibility of both weights,
    \[
      \E[D_\infty\mid w^m, w^M\in W_{\mathrm{adm}}] = (1/m - 1/M)\cdot\sum_k\bar w_k\tr\bigl(\tilde S_k\tilde C\tilde S_k^\T\bigr)
      + O(m^{-3/2}) + O(m^{-1}M^{-1/2}):
    \]
    the $1/m$ law with an explicit constant, and the $(1/m - 1/M)$ law of the shared-particle coupling observed in
    Theorem D$'$'s numerics.
  \item \emph{(the fluctuations)} As $M\to\infty$ ($w^M\to\bar w$) and then $m\to\infty$,
    $\sqrt m(w^m - \bar w)\to\mathcal{N}(0,\tilde C)$ and $\sqrt m\cdot(\vartheta_i - \bar\vartheta_i)\to S_{k(i)}\xi$
    with $\xi\sim\mathcal{N}(0,\tilde C)$ for every particle $i$ of cluster $k(i)$: the fluctuations of the coupling
    are Gaussian, carried entirely by the cluster weights, and all particles of a cluster are displaced by the same
    vector.
  \item \emph{(admissibility)} Let
    $\delta_w := \sup\{\delta : \norm{w - \bar w}_\infty\le\delta\Rightarrow w\in W_{\mathrm{adm}}\} > 0$. By
    Hoeffding's inequality for each multinomial coordinate,
    $\mathbb{P}(w^m\notin W_{\mathrm{adm}})\le2K\exp(-2m\delta_w^2)$: the finite-$m$ configuration is a strict local
    minimiser with probability exponentially close to one, the exceptional event being the one on which Remark
    D$'''$.2's floor appears on a warm start --- a cluster of the $m$-system splits and particles migrate to other
    basins, because the realised weights make the replicated configuration a saddle.
\end{enumerate}
\end{thm}
\status{Formal: (i) and the algebra of (ii)--(iv) are exact (the block splitting is Proposition K(i) cluster by
cluster, confirmed on $J_{10}$ to $8\cdot10^{-12}$ in referee round 9; implicit function theorem; the multinomial
covariance of shared labels; Hoeffding), the convergence of both systems to the stated configurations is assumed (it
is what Theorem D$'''$ proves for $m\ge m_0$ and what the numerics show for $m = 16$--$256$ at admissible weights; at
inadmissible weights the $m$-system sits on the saddle for a long time and then leaves it, numerics below). Checked:
\code{weighted_stationary_atoms}, \code{weight_sensitivity}, \code{predicted_stationary_D} in
\code{chaos_learning.py}; on the instance below the replicated weighted stationary point is a critical point of
$J_{16}$ to $10^{-10}$, the linearisation $\vartheta^*(\bar w + \mathrm{d}w) - \vartheta^* = S\cdot\mathrm{d}w$ has a
second-order remainder, and the per-cluster displacements of the coupled run point along $S_k(w^m - w^M)$ (test and
table below). The proposition is the sharp form of Theorem D$'''$'s constant:
$S = -H_w^{-1}\partial_w\nabla\mathcal{G}_w$ has the structure (stiffness)$^{-1}$ $\times$ (feature correlations), of
which $C_1 = 4e^{2\Lambda_1}(L_\vartheta^G)^2K_\Phi^2/\kappa^2$ is the worst-case envelope ($\kappa^{-1}$ for the
stiffness, $L_\vartheta^GK_\Phi$ for the correlations).}

\begin{rmk}{Remark}{D$'''$.6}[the regime of the theorem, and what the numerics test]
The sample size $m_0 = 16K_\Phi^2/\theta^2$ from which the theorem's tail is in force is astronomically large on
realistic instances ($10^{23}$ below): $\theta$ carries $\varrho'\kappa^2e^{-2\Lambda_1}/(L_\vartheta^G)^3$ from the
per-particle bootstrap of Step~6, because the interaction's effect on one particle is bounded by
$(L_\vartheta^G)^2D^{1/2}$ and divided by the stiffness $\kappa$. This is not an artefact of the proof's constants
only: Proposition D$'''$.5 says the particles of a cluster are displaced by $S_k(w^m - w^M)$, of order
$\norm{S}m^{-1/2}$ with $\norm{S}\sim(\text{stiffness})^{-1}$, which exceeds any fixed $\varrho'$ for $m$ below
$(\norm{S}/\varrho')^2$ --- on the instance below the $m = 16$ particles sit $0.4$--$1.4$ whitened units from their
reference copies, outside the neighbourhood $N_{\varrho'}$ in which $\mathcal{V}[\mu_{\bar\Theta}]$ is convex, at the replicated stationary configuration $\Theta^*_{16}(w^m)$ of
their own $J_{16}$ --- a strict local minimiser when $w^m\in W_{\mathrm{adm}}$ and a saddle otherwise (Proposition
D$'''$.5(i), (iv); one of the four $m = 16$ seeds below). Hypothesis (H) is nevertheless the right object: it is the condition
under which the reference configuration is stable to splitting, and the $m$-system's configuration is the same
configuration with perturbed weights. What the numerics test is therefore (a) hypothesis (H) at the stationary
measure and along the flow, (b) the mechanism --- the coupled distance saturates (no growth) and follows the $1/m$
law --- and (c) the sharp constant of Proposition D$'''$.5; the theorem's own constant $C_1$ is reported for the
record.
\end{rmk}

\begin{thm}{Corollary}{D$'''$.7}[the real loop on a warm start: uniform in time through Theorem J]
Let $\mu^* = \sum_k\bar w_k\delta_{\vartheta^*_k}$ be a $K$-atom stationary measure as in Proposition D$'''$.5,
satisfying the splitting condition $\sigma_0\Ginv - R[\mu^*](\vartheta^*_k)\succ0$ at every atom and the
non-degeneracy of the $K$-atom Hessian, and let $w\mapsto\Theta^*(w)$ be the weighted stationary point of
D$'''$.5(i). Fix $m\le M$ and cluster weights $w^m, w^M\in W_{\mathrm{adm}}$ (Proposition D$'''$.5(i)), and let
$\Theta^*_m :=$ the configuration with $n_k = mw^m_k$ copies of $\vartheta^*_k(w^m)$ (likewise $\Theta^*_M$). Both are
non-degenerate local minimisers of $J_m$ and $J_M$ (Proposition D$'''$.5(i)), hence hyperbolic critical points, and
Theorem J applies to each loop: there are $k_0$, $\rho_0$, $C$ (depending on $\Theta^*_m$ and $\Theta^*_M$, hence on
$m$, $M$ and the weights) such that for $k_e\ge k_0$ the real closed loop \textup{(CL)} with $m$ kernels started at
$\norm{\Theta(0) - \Theta^*_m}_{\Ginv}\le\rho_0$, $\norm{e(0)}\le R$ converges to its $T_0$-periodic orbit
$\Theta_p^m$, and likewise the $M$-kernel loop to $\Theta_p^M$. Under the shared-particle coupling (the $m$ kernels
are the first $m$ of the $M$, the plants are driven by the same reference) the coupled distance
$D(t) := \frac1m\sum_i\norm{\vartheta_i(t) - \bar\vartheta_i(t)}^2_{\Ginv}$ satisfies, for all $t\ge0$,
\begin{multline*}
  D(t)^{1/2}\le D_{\mathrm{orb}}^{1/2} + Ce^{-\mu t}\bigl(\norm{\Theta(0) - \Theta_p^m(0)}_{\Ginv} + \norm{x(0) - x_p^m(0)}\bigr)\\
  + Ce^{-\mu t}\bigl(\norm{\bar\Theta(0) - \Theta_p^M(0)}_{\Ginv} + \norm{\bar x(0) - x_p^M(0)}\bigr),
\end{multline*}
with $\mu = (\lambda_*/k_e)(1 - \kappa_J)$ the rate of Theorem J(c) ($\kappa_J$ its small parameter; the smaller of
the two loops' rates) and
\begin{align*}
  D_{\mathrm{orb}} &:= \sup_t\frac1m\sum_i\norm{\Theta_p^m(t)_i - \Theta_p^M(t)_i}^2_{\Ginv}\\
  &= \sum_kw^m_k\norm{\vartheta^*_k(w^m) - \vartheta^*_k(w^M)}^2_{\Ginv} + O(1/k_e)\\
  &= \sum_kw^m_k\norm{S_k(w^m - w^M)}^2_{\Ginv} + O(\norm{w^m - w^M}^3) + O(1/k_e):
\end{align*}
the coupled distance of the real loop is bounded uniformly in time, converges to the distance between the two
periodic orbits, and that distance is Proposition D$'''$.5's stationary distance up to the $O(1/k_e)$ offsets of
Theorem J(a),(d) (which differ between the two loops by $O(k_e^{-1}\norm{w^m - w^M}) + O(k_e^{-2})$, the first-order
offset of a replicated configuration being a smooth function of the $K$-atom data $(\Theta^*(w),w)$ by permutation
equivariance). In particular, for cluster labels i.i.d.\ of law $\bar w$ and conditionally on both weights being
admissible and both starts lying in their basins,
$\E D_{\mathrm{orb}} = (1/m - 1/M)\sum_k\bar w_k\tr(\tilde S_k\tilde C\tilde S_k^\T) + O(m^{-3/2}) + O(1/k_e)$.
\end{thm}
\status{Proved as stated (Theorem J(a),(c),(d) applied to each loop, the triangle inequality in the $\Ginv$ metric,
and Proposition D$'''$.5(ii) for the distance between the critical points), under Theorem J's hypotheses for both
configurations: the constants $k_0$, $\rho_0$, $C$, $\mu$ are those of Theorem J at $\Theta^*_m$ and $\Theta^*_M$,
and the corollary does not assert their uniformity in $m$ or $M$ (Remark D$'''$.8). It is the statement of Remark
D$'''$.4 for a warm start: the real loop's coupled distance saturates, at a value set by the cluster weights.
Consistent with the instance of this section's numerics (table 5: the real loop at $k_e = 25$ from the jittered
mixture over 80 slow units, a sixth of the relaxation time $1/\kappa$, follows its averaged slow dynamics on the same
initial clouds to within a few percent at every window end; the saturation at the orbit distance lies beyond the
horizon run and is not checked).}

\begin{proof}
Theorem J(c) for the $m$-kernel loop gives
$\norm{\Theta(t) - \Theta_p^m(t)}_{\Ginv}\le Ce^{-\mu t}(\norm{\Theta(0) - \Theta_p^m(0)}_{\Ginv} + \norm{x(0) - x_p^m(0)})$
and the same for the $M$-kernel loop; the first $m$ rows of the $M$-loop are the reference particles
$\bar\vartheta_i$. For the $(m\times p)$-row norm
$\frac1m\sum_i\norm{\cdot}^2_{\Ginv}\le\norm{\cdot}^2_{\Ginv}/m\le\norm{\cdot}^2_{\Ginv}$ (the full configuration
norm dominates the mean square over the first $m$ rows), so by the triangle inequality
$D(t)^{1/2}\le[\frac1m\sum_i\norm{\Theta_p^m(t)_i - \Theta_p^M(t)_i}^2]^{1/2} + (\text{the two transients})$.
Theorem J(a) places each orbit within $C/k_e$ of its critical point, uniformly in $t$, so the first term is
$\sum_kw^m_k\norm{\vartheta^*_k(w^m) - \vartheta^*_k(w^M)}^2_{\Ginv}$ up to $O(1/k_e)$ (the two offsets are
$O(1/k_e)$ each, with $O(k_e^{-2})$ remainders; at a replicated configuration the first-order offset is the same
vector for all particles of a cluster and depends smoothly on the $K$-atom data $(\Theta^*(w),w)$, by the permutation
equivariance of \textup{(CL)} and the block splitting of Proposition D$'''$.5(i), so the two loops' offsets differ by
$O(k_e^{-1}\norm{w^m - w^M}) + O(k_e^{-2})$ --- the two loops live on $\R^{mp}$ and $\R^{Mp}$, and it is this
reduction to the $K$-atom data that makes them comparable), and Proposition D$'''$.5(ii) expands the distance between
the critical points in $w^m - w^M$. The expectation statement is D$'''$.5(ii)'s.
\end{proof}

\begin{rmk}{Remark}{D$'''$.8}[what the corollary does and does not say]
It transfers the uniform-in-time statement to the real loop in the regime where Theorem J applies to both loops ---
a warm start in the basins of two hyperbolic minimisers at a gain $k_e$ above both thresholds. Theorem J's constants
at a replicated configuration are controlled by the $K$-atom data (the Hessian of $J_m$ there is the collective block
$\nabla^2\mathcal{G}_{w^m}\oplus{}$the separation blocks $\sigma_0\Ginv - R_k$ of Proposition K(i), both independent
of $m$ beyond the weights), so the threshold and the basin are expected to be uniform in $m$ for configurations with
the same weights; the corollary does not prove this, and the cold start --- where the coupled distance of the real
loop grew in Theorem D$''$'s experiment --- is outside it on the same two counts as in Remark D$'''$.2.
\end{rmk}

% Numerics (PARADIGM.md section 5d, third pass; the numbers are PARADIGM's copies of results/chaos/learning_table.md).
\paragraph{Numerics} (\code{exp_chaos_learning.py} $\to$ \code{results/chaos/learning_table.md}). The instance:
$n = 1$, saturated features with $a_{\max} = 4$, $c_{\max} = 3$, $s\in[0.3,2]$, prior $\vartheta^0 = 0$,
$G = \operatorname{diag}(1,0.25,0.25)$; a plant of three Gaussian bumps, $(A,C,S) = (1,-1.5,0.5)$, $(-0.7,0,0.4)$,
$(0.8,1.5,0.5)$, which three kernels of weights $3A_k$ represent exactly; reference $x_m = 2\sin t$, $\rho$ on a grid
of 256 times; slow time, midpoint rule with step $0.02$; $m\in\{16,64,256\}$ sharing their particles with an
$M = 2048$ reference (the stand-in for the mean-field flow, as in Theorems D$'$ and D$''$); horizon 900 slow units;
four seeds per start. Proposition K gives $\sigma_{\mathrm{crit}} = 0.04716$ (collapsed atom $a = 0.331$,
$c = -0.057$, $s = 1.218$); the runs use $\sigma_0 = 0.15\sigma_{\mathrm{crit}} = 0.00707$ (leak time 141 slow units).
Whitened constants by search: $L_\vartheta^G = 4.044$ (closed form $4.658$), $L_2^G = 11.11$ (closed form
$\lambda_Ga_{\max}/s_{\min}^2 = 44.44$); $K_\Phi = 4$, $F_\rho = 0.999$.

\smallskip
\noindent\emph{The hypothesis at the stationary measure, and the window (table 1).} The critical point of $J_3$
reached from the exact representation, atoms
\[
  (a,c,s) = (2.4504,-1.4161,0.8287),\quad(-2.2879,0.0120,0.7009),\quad(2.0478,1.3427,0.8810),
\]
has $\norm{\nabla J_3}_\infty = 4\cdot10^{-17}$ (a stationary measure $\mu^*$ of the
mean-field flow), fit error $\norm{f - \fhat}/\norm{f} = 0.244$ on the orbit (the leak's bias: amplitudes
$2.0$--$2.5$ against $3A_k = 2.1$--$3$, widths $0.70$--$0.88$ against $0.4$--$0.5$), and a positive whitened Hessian
(eigenvalues $0.0057$--$0.2003$). $\kappa_R$ at the atoms is $0.00004$, $-0.00227$, $0.00411$; over the spheres of
radius $\varrho'$ around them (400 directions, radii $\varrho'$ and $\varrho'/2$) it is $0.00461$
($\varrho' = 0.05$), $0.00512$ ($0.1$), $0.00619$ ($0.2$) and $0.00730$ ($0.3$), so (H) holds at $\mu^*$ with
$\varrho' = 0.1$ and margin $\kappa = \sigma_0 - 0.00512 = 0.00195$ (relaxation time $1/\kappa = 512$) and fails from
$\varrho'\approx0.3$ on. The window of Remark D$'''$.1 is realised:
$\kappa_R[\mu^*](N_{0.1}) = 0.00512 < \sigma_0 = 0.00707 < \sigma_{\mathrm{crit}} = 0.04716$. At
$\sigma_0 = 0.3\sigma_{\mathrm{crit}}$, by contrast, the three-atom critical point (fit error $0.406$) fails the
splitting condition at its third atom ($\kappa_R = 0.01696 > \sigma_0 = 0.01415$): there the mean-field limit has
more atoms, and the hypothesis has to be checked at that configuration. The theorem's constants for the warm start
($\Lambda_1 = 0$, $\kappa := \frac12\min\{\kappa(\varrho'),\min_sa(s)\} = 0.00070$, the factor $\frac12$ leaving room for the
sampling of the supremum by finitely many directions): $R_\Theta = 2857$,
$\varrho = 1.1\cdot10^{-6}$, $\theta = 4.6\cdot10^{-11}$, $C_1 = 2.2\cdot10^9$, $m_0 = 1.2\cdot10^{23}$ --- the
worst-case envelope of Remark D$'''$.6.

\smallskip
\noindent\emph{The hypothesis along the flow (table 2).} On the warm start ($\mu_0$ the mixture of Gaussians of
width $0.02$ around the atoms) the margin $a(s) = \sigma_0 - \kappa_R[\mu^*]($spheres of radius $0.1$ around 256 reference
particles, 24 directions$)$ stays between $+0.00139$ and $+0.00167$ over the whole run, so (H) holds
with $s_1 = 0$; the same statistic with the reference's own residual is larger by at most $0.00013$ (its sampling
fluctuation; the proof's bound (0.2) allows $L_2^G\norm{r_{\mathrm{ref}} - r^*}_{L^1(\rho)}\le0.214$), the
reference's rms drift decays from $6\cdot10^{-3}$ to $3\cdot10^{-6}$ and its fit error stays at $0.245$--$0.248$: the reference is stationary. On the cold start ($a\sim\mathcal{N}(0,0.8^2)$, $c\sim\mathcal{U}(-2,2)$,
$s\sim\mathcal{U}(0.4,1)$, a connected cloud) the margin is negative over the whole run, $-0.196$ at the first record ($-0.179$ at $s = 1$) and $-0.0025$ at
$s = 900$, while the fit error falls from $0.948$ to $0.244$ --- the same fit as the warm limit --- and
the rms drift from $0.15$ to $4.7\cdot10^{-4}$: the cold reference converges to a configuration of the same quality
whose support reaches the ridges (Remark D$'''$.2).

\smallskip
\noindent\emph{Warm start (table 3).} Window maxima of $D_m$ over $(0,112.5]$, $(112.5,225]$, $(225,450]$,
$(450,900]$ (Table~\ref{tab:thmDppp:warm}):

\medskip
\noindent\begin{minipage}{\textwidth}
\centering\footnotesize
\captionof{table}[Theorem D$'''$: the warm start]{Theorem D$'''$ on the three-bump instance, warm start: window
maxima of the coupled distance $D_m$ over four windows of slow time, the mean on the last window, the theorem's
constant $C_1/m$, the learned-function distance and the per-particle maximum
(\code{results/chaos/learning_table.md}, table 3).}
\label{tab:thmDppp:warm}
\setlength{\tabcolsep}{3pt}
\begin{tabularx}{\textwidth}{@{}c>{\centering\arraybackslash}X>{\centering\arraybackslash}X>{\centering\arraybackslash}X>{\centering\arraybackslash}X>{\centering\arraybackslash}X>{\hsize=0.7\hsize\centering\arraybackslash}X>{\hsize=1.2\hsize\centering\arraybackslash}X>{\hsize=1.1\hsize\centering\arraybackslash}X@{}}
\toprule
$m$ & $(0,112.5]$ & $(112.5,225]$ & $(225,450]$ & $(450,900]$ & mean on $(450,900]$ & $C_1/m$
  & mean $\norm{\fhat_\Theta - \fhat_{\mathrm{ref}}}$, last window & $\max_i\norm{\delta_i}_{\Ginv}$ at $s = 900$\\
\midrule
16 & 0.396 $\pm$ 0.24 & 0.423 $\pm$ 0.26 & 0.441 $\pm$ 0.28 & 0.448 $\pm$ 0.28 & 0.446 $\pm$ 0.28 & $1.4\cdot10^8$
  & 0.0307 $\pm$ 0.0097 & 1.36\\
64 & 0.122 $\pm$ 0.046 & 0.131 $\pm$ 0.049 & 0.136 $\pm$ 0.050 & 0.138 $\pm$ 0.050 & 0.138 $\pm$ 0.051 & $3.4\cdot10^7$
  & 0.0154 $\pm$ 0.0033 & 0.62\\
256 & 0.0119 $\pm$ 0.0050 & 0.0127 $\pm$ 0.0052 & 0.0134 $\pm$ 0.0055 & 0.0137 $\pm$ 0.0056 & 0.0136 $\pm$ 0.0055
  & $8.4\cdot10^6$ & 0.00467 $\pm$ 0.0012 & 0.21\\
\bottomrule
\end{tabularx}
\end{minipage}
\medskip

The distance saturates: last-to-first window ratios $1.13$, $1.13$, $1.15$ and last-to-previous $1.02$, $1.01$,
$1.02$. Its $m$-law has slope $-1.28\pm0.30$ (per-seed fits; the theorem's mean square: $-1$), the learned function's
$-0.69\pm0.16$ ($m^{-1/2}$: $-0.5$). The per-particle maxima $1.36$, $0.62$, $0.21$ are far above
$3\varrho'/4 = 0.075$: the localisation of part (c) is not in force at these $m$ (Remark D$'''$.6), and the values
lie below the vacuous $C_1/m$ by 8.4--8.8 orders of magnitude.

\smallskip
\noindent\emph{Proposition D$'''$.5 (the same table).} The sensitivities of the three atoms to the weights have
whitened norms $\norm{S_k} = 8.38$, $4.66$, $8.32$. The multinomial expectation
$\E D_\infty = (1/m - 1/M)\sum_k\bar w_k\tr(\tilde S_k\tilde C\tilde S_k^\T)$ is $0.823$, $0.201$, $0.0454$ against
the four-seed means $0.446\pm0.28$, $0.138\pm0.051$, $0.0136\pm0.0055$ (ratios $0.54$, $0.69$, $0.30$: a mean over
four seeds of a heavy-tailed quantity, $\E\abs{w^m - w^M}^2$ being carried by the rare large fluctuations). The
per-seed law is the sharp test: $\sum_kw^m_k\norm{\tilde S_k(w^m - w^M)}^2$ computed from the realised cluster
weights predicts the measured $D_m$ at $s = 900$ within 25\% in 9 of the 12 (seed, $m$) pairs (for instance
$0.123\to0.112$, $0.187\to0.186$, $0.206\to0.205$ at $m = 16$; $0.0303\to0.0261$, $0.00087\to0.00085$,
$0.0081\to0.0087$, $0.0187\to0.0190$ at $m = 256$); the largest deviation, a ratio $0.57$ at $m = 16$, is the seed
with the largest weight fluctuation $\abs{w^m - w^M} = 0.275$, where the linearisation's remainder and the saturation
$a\le a_{\max}$ enter. The exact first equality of (ii) --- $\sum_kw^m_k\norm{\vartheta^*_k(w^m) - \vartheta^*_k(w^M)}^2$
with the weighted stationary points themselves --- reproduces the measured $D_m$ within 2.3\% in all 12 pairs
($1.292\to1.292$ for the seed just mentioned): every $m$-system sits at the replicated weighted stationary
configuration. The splitting margin of $\Theta^*(w^m)$ is positive in 11 of the 12 pairs ($+0.0006$ to $+0.0050$);
it is negative, $-0.0022$, for $m = 16$, seed 3, whose weights $(9,4,3)/16$ lie outside the admissible set: there
the replicated configuration is a saddle of $J_{16}$ and the plateau at $s = 900$ is the saddle's --- referee round 9
continued that run to $s = 3600$ and saw $D$ grow from $1.29$ to $5.13$ as two particles changed basins and the
weights re-sorted to $(7,5,4)/16$ (Proposition D$'''$.5(iv); the root-mean-square weight fluctuation is $0.204$,
$0.102$, $0.051$ at $m = 16$, $64$, $256$ against admissible margins of $0.06$--$0.36$ in $\norm{\cdot}_\infty$). Per
cluster (seed 0), the measured mean displacements of the $m$-system's particles point along $S_k(w^m - w^M)$ with
cosines $0.92$--$1.000$ and norms within 0.8--20\% ($m = 16$: $0.434$ against $0.455$, $0.130$ against $0.129$,
$0.338$ against $0.358$): the particles of a cluster move together, by the predicted vector. The ratios $0.54$,
$0.69$, $0.30$ of the four-seed means to the multinomial expectation are the sampling of a heavy-tailed quantity at
$m = 16$ and $64$ and, at $m = 256$, four seeds whose realised fluctuations happen to be small.

\smallskip
\noindent\emph{Cold start (table 4).} Window maxima of $D_m$: $1.81\pm0.28\to3.13\to4.50\to5.33\pm0.16$
($m = 16$), $0.505\to0.900\to1.42\to1.79\pm0.93$ ($64$), $0.0744\to0.162\to0.347\to0.583\pm0.19$ ($256$): the
distance grows over the whole horizon (last-to-first $2.95$, $3.54$, $7.84$; last-to-previous $1.19$, $1.26$,
$1.68$), with slope $-0.87\pm0.14$ in $m$ (the learned function $-0.66\pm0.06$, steeper than the $-0.5$ of the warm start). The
per-particle deviations are
bimodal: medians $0.40$, $0.14$, $0.067$ (the $m^{-1/2}$ law of the particles that stayed with their references)
against maxima $6.7$, $6.3$, $6.2$ for every $m$ (the separation between basins), and the fraction of pairs more than
one whitened unit apart is $0.203\pm0.047$, $0.082\pm0.021$, $0.037\pm0.010$ --- a factor $2.5$ and $2.2$ per factor 4 in $m$, a decay like $m^{-0.64\pm0.08}$
(per-seed fits; the medians decay like $m^{-0.62\pm0.07}$) --- steeper than the $m^{-1/2}$ of Remark D$'''$.2's
boundary-layer count, which is a lower-order prediction (the flips contribute $p_m\cdot O(R^2)\approx0.2\cdot38 = 7.8$,
$3.1$, $1.4$ against the measured $5.3$, $1.8$, $0.58$). This is Theorem
D$''$'s negative finding reproduced and diagnosed on the averaged system: a connected cold start that splits between
basins is outside (H), and its coupling distance is dominated by the pairs that chose different basins.

\smallskip
\noindent\emph{The real loop (table 5).} The closed loop at $k_e = 25$ ($\sigma = \sigma_0/k_e = 2.8\cdot10^{-4}$,
$\tau = 0$, $\mathrm{d}t = 10^{-3}$, three seeds) from the same warm start: window maxima over $(0,500]$,
$(500,1000]$, $(1000,2000]$ real-time units are $0.328\pm0.11\to0.467\to0.590\pm0.17$ ($m = 16$),
$0.0605\to0.0801\to0.0909\pm0.045$ ($64$), $0.0188\to0.0253\to0.0294\pm0.011$ ($256$), slope $-1.13\pm0.15$ in $m$;
the reference's fit error stays at $0.245$--$0.246$. The horizon is 80 slow units, a sixth of the relaxation time
$1/\kappa = 512$, and the growth (last-to-first $1.80$, $1.50$, $1.56$) is the early transient of the averaged dynamics: the averaged
flow started from the real loop's own initial clouds (reproduced from their seeds) gives
$D_{\mathrm{real}}(t)/D_{\mathrm{averaged}}(t/k_e)$ between $0.993$ and $1.004$ at every window end, for every seed
and every $m$ --- the real loop at $k_e = 25$ follows its averaged slow dynamics to within 0.7\% (across different
clouds the last-window maxima are $1.57$, $0.78$ and $2.58$ times the averaged runs' $0.377$, $0.117$, $0.0114$,
within the seed-to-seed spread). Two of the three $m = 16$ clouds have inadmissible or marginal weights (splitting
margins $-0.0021$, $-0.00001$; the third $+0.0026$), so their eventual configuration is a saddle's; at $m = 64$ and
$256$ all margins are positive ($+0.0029$ to $+0.0049$). The saturation at the orbit distance of Corollary
D$'''$.7 lies beyond the horizon run.

\smallskip
\noindent\emph{Reading.} In the window $\kappa_R[\mu^*] < \sigma_0 < \sigma_{\mathrm{crit}}$ the network learns
(the collapsed configuration is unstable, the three kernels fit the plant to $0.244$) and propagation of chaos is
uniform in time on a warm start with admissible weights: the coupled distance saturates with the $1/m$ law at the
stationary distance of the replicated weighted configuration, which the exact form of Proposition D$'''$.5(ii) gives
within 2.3\% --- the level is the weight-sensitivity law, not the theorem's constant. At $m = 16$ the multinomial
weights leave the admissible set with appreciable probability; the one such seed sits on a saddle and later splits a
cluster. The cold start of a connected cloud shows what the hypothesis excludes: growth over the horizon and a floor
of pairs in different basins whose fraction decays like $m^{-0.64}$, while the learned function still converges at a
rate between $m^{-1/2}$ and $m^{-2/3}$ in both cases. The real loop at $k_e = 25$ follows its averaged slow dynamics
to within 0.7\% over the horizon run.

\checkedby{\code{tests/test_chaos_learning.py}: \textup{($\mathfrak{A}_m$)}'s drift equals $-G\nabla J_m$ of
\code{periodic.ClosedLoop} to machine precision (\code{test_averaged_slow_loop_is_the_G_gradient_flow_of_J_m}); the
three-term identity of Step~1 holds to rounding on random configurations with a mean-field stand-in, the interaction
term equals $-\norm{\fhat_\Theta - \fhat_{\bar\Theta}}^2$ up to a relative $O(D^{1/2})$ correction (ratio
$0.98\to0.9999$ as $\delta$ is scaled by $1\to0.01$) and obeys (3.1) at every scale\newline
(\code{test_coupling_identity_three_terms_and_the_sign_of_the_interaction}); $\kappa_R$ at the atoms of a collapsed
configuration with the collapsed residual equals Proposition K's $\lambda_{\max}(G^{1/2}R_2G^{1/2})$, which equals
$\sigma_0$ at $\sigma_0 = \sigma_{\mathrm{crit}}$, the collapsed critical point has zero drift, and the deterministic
neighbourhood search over the sphere of radius $\varrho$ reproduces the first-order expansion
$\lambda_{\max} + \varrho\norm{G^{1/2}\nabla\lambda_{\max}}$ with a relative error $O(\varrho)$
(\code{test_residual_curvature_}\dots); the searched whitened constants reproduce the document's values and lie below
their closed forms, \code{theorem_constants} satisfies the four constraints the bootstrap needs, and $\Lambda_1$
integrates exactly to $s_1$ (\code{test_whitened_constants_}\dots); at $\sigma_0 = 0.15\sigma_{\mathrm{crit}}$ the
three-atom critical point of $J_3$ has zero drift and (H) holds at it with $\varrho' = 0.1$ inside the learning
regime, the replicated weighted stationary point is a critical point of $J_{16}$ to $10^{-10}$, the linearisation in
the weights has a second-order remainder, and on a coupled run from the jittered mixture ($m = 16$, $64$ against
$M = 512$, 60 slow units) $\kappa_R[\mu^*]$ over the neighbourhood of the reference's support stays below $\sigma_0$,
the two curvature statistics differ by at most the bound ($0.2$), the per-cluster displacements point along
$S_k(w^m - w^M)$ (cosine $\ge0.8$) with $D$ within a factor 2 of $\sum_kw^m_k\norm{S_k(w^m - w^M)}^2$,
$\norm{\fhat_\Theta - \fhat_{\bar\Theta}}\le L_\vartheta^G\cdot\operatorname{mean}\norm{\delta_i}$, $D_{64} < D_{16}$
and the (vacuous) bound $C_1/m$ holds\newline
(\code{test_hypothesis_H_at_the_three_atom_stationary_measure_and_the_weight_sensitivity_law}).}
